\section{问题分析}\label{sec:problem-analysis}

三个问题共享机场、海上设施、机型、燃油及停靠等基本规则，但约束逐步增加：问题一处理单向出海运输，问题二扩展为多方向联合取送，问题三进一步加入时间窗、具体机队和多日排班。因此，三问可视为同一运输优化框架下逐层扩展的三个阶段。从车辆路径问题的视角看，异构车队、多中心与多架次均属于现实运输优化中常见的扩展结构，已有研究分别讨论了异构车队多中心协同配送以及异构车队多架次路径优化问题\cite{liang2024,sethanan2020}。

\subsection{问题一分析}\label{subsec:q1-analysis}

问题一包含1600名出海人员，运输方向均为陆地机场至海上设施。在没有海上上客的情况下，同一架次中的人员只会随停靠过程逐步下机，因此机上人数单调减少，最大载客量出现在起飞后的第一航段，这使问题一的容量关系相对简单。

该问题的关键在于，在满足机型容量、最多海上着陆次数、燃油安全及机场闭环等条件下，合理组合人员并确定海上设施访问顺序。不同路线不仅影响飞机总使用时间，也会影响人员在途时间、燃油消耗及座位利用率。因此，可将问题转化为可行飞行路线构造与人员分配相结合的运输优化问题，以总飞机使用时间为主要优化目标。

\subsection{问题二分析}\label{subsec:q2-analysis}

问题二将运输规模扩大至4000人，并新增海返人员和海上设施间穿梭人员。与问题一相比，主要变化不是人数增加，而是人员可以在飞行途中不同站点上下机，使飞机载客量随航段动态变化。

在同一设施完成下客后，释放出的座位可以立即用于后续人员上机，因此需要按照路线访问顺序逐航段计算实际载客人数。同时，海返和穿梭人员只有在飞机先访问其起点、再访问终点时才能完成运输，人员分配与路线顺序之间形成更强耦合。

因此，问题二应在问题一路线规划基础上进一步考虑人员取送顺序和动态座位复用，形成路线选择、人员上下机安排与航段容量联合优化的问题。

\subsection{问题三分析}\label{subsec:q3-analysis}

问题三在问题二的联合运输基础上进一步加入7日规划周期、24架具体直升机以及人员时间窗和任务类别，其中3840名非临时人员必须完成运输，160名临时人员允许选择性接纳。此时，一条空间上可行的运输路线还必须确定具体执行日期、起飞时刻和执行飞机，并满足人员最早出发、最晚到达以及飞机周转和每日运营时间等要求。

因此，问题三由前两问的路线与人员配置问题进一步扩展为人员、路线、时间与具体飞机的联合排班问题。已有海上油气人员直升机调度研究将同类问题描述为同时具有取送、异构机队、多基地、多架次和时间同步特征的复杂车辆路径问题\cite{nafstad2021}，这与本题第三问的主要约束结构具有较强对应性。同时，由于非临时任务必须优先保障，而临时任务允许选择性接纳，可先完成非临时任务并建立飞机使用时间基准，再在不超过该基准的条件下尽可能接纳临时人员；在固定临时人员接纳数量后，进一步按照既定优先级优化总飞机使用时间、人员在途时间、燃油消耗、架次数及座位利用率。

三个问题的核心递进关系可概括为
\[
\boxed{
\text{单向运输}
\rightarrow
\text{联合取送与动态容量}
\rightarrow
\text{时间窗约束下的有限机队排班}
}.
\]
